\subsection{RMRにおける冗長性評価}
本実装では，関連研究のRMRの考え方に基づき，観測頻度，動的変化量，距離を用いて物標情報の冗長性を評価する．
物標$i$の冗長性スコア$R_i$を次式で定義する．
\begin{equation}
 R_i =
 \frac{N_i}
 {D_i \min(\Delta P_i, \Delta V_i)}
\end{equation}
ここで，$N_i$は一定時間窓内における物標$i$の観測回数，$D_i$は自車両から物標$i$までの距離，$\Delta P_i$は前回CPMに含めた時点からの位置変化量，$\Delta V_i$は速度変化量である．
$R_i$が大きい物標ほど，頻繁に観測され，動的変化が小さく，周辺車両でも認識しやすい冗長な情報とみなし，CBRに応じた削除個数の範囲でCPMから削除する．

\subsection{MEC/V2N2V経路の補足パラメータ}
本評価では，計算負荷の削減と評価対象の切り分けのため，MEC経路を詳細なセルラ無線モデルではなく理想リンクとして扱った．
教授への補足説明で重要なパラメータを表\ref{tab:mec_params}に示す．

\begin{table}[t]
  \centering
  \caption{MEC/V2N2V経路の主な補足パラメータ}
  \label{tab:mec_params}
  \begin{tabular}{ccc}
    \hline
    項目 & 設定値 & 意味 \\ \hline \hline
    MEC理想リンク & 有効 & セルラ網詳細モデルを省略 \\
    転送範囲 & 200\,m & MECから周辺車両へ再配信する範囲 \\
    下りPDR & 0.97 & MECから車両への到達率 \\
    遅延範囲 & 20--100\,ms & MEC経路遅延の範囲 \\
    代表遅延 & 50\,ms & MEC経路の平均的な遅延 \\
    低重要度最大転送確率 & 0.5 & 低重要度情報をV2N2Vにも送る最大確率 \\ \hline
  \end{tabular}
\end{table}

V2N2V送信量は，車両からMECへのアップリンク量と，MECから周辺車両への再配信量の合計である．
そのため，MECから複数車両へ転送される分，V2N2V送信量はV2V送信量より大きくなる．

\subsection{背景トラフィック条件}
本評価では，CPM以外の通信負荷を模擬するため，各車両にNR sidelink背景トラフィックを発生させた．
背景負荷は道路長方向の位置に応じて変化し，幅500\,mの高負荷帯が道路上を移動するように設定した．
判定には車両の$x$座標のみを用いるため，同じ道路長方向位置にいる車両であれば，対向車線の車両も同じ背景負荷を受ける．

\begin{table}[t]
  \centering
  \caption{移動背景トラフィックの主な設定}
  \label{tab:bg_params}
  \begin{tabular}{ccc}
    \hline
    項目 & 設定値 & 意味 \\ \hline \hline
    BGパケットサイズ & 800\,byte & 基準となる背景通信サイズ \\
    BG送信間隔 & 80\,ms & 各車両の背景通信周期 \\
    高負荷帯幅 & 500\,m & 中心から前後250\,mの帯状領域 \\
    高負荷帯速度 & 25\,m/s & 道路$x$方向への移動速度 \\
    高負荷係数 & 1.8 & 高負荷帯内のパケットサイズ倍率 \\
    最小負荷係数 & 0.2 & 高負荷帯外のパケットサイズ倍率 \\ \hline
  \end{tabular}
\end{table}

高負荷係数は，設定値 \texttt{NR\_BG\_WAVE\_AMPLITUDE=0.8} により $1.0+0.8=1.8$ として与えられる．
したがって，高負荷帯内ではBGパケットサイズが $800\times1.8=1440$\,byte，高負荷帯外では $800\times0.2=160$\,byte となる．
これらの値はCBRを直接指定するものではなく，背景トラフィック量を変化させるための係数であり，実際のCBRはCPM通信量や車両密度を含む通信状況から結果として決まる．
